The second example demonstrates the pipeline's ethical safeguard: when the
context-extraction module identifies a conversation as emotionally sensitive,
the system suppresses ad delivery entirely and returns a purely supportive
response.  This gate operates independently of the ad inventory --- even if
highly relevant ads existed, they would not be shown.

\begin{figure}[H]
    \centering
    \includegraphics[width=\textwidth]{images/illustrative_example_2.png}
    \caption{Chat response to a sensitive mental-health query.  No sponsored
    content appears; the assistant responds with empathetic, helpful guidance
    entirely free of product mentions.}
    \label{fig:example-2}
\end{figure}

\subsubsection{Stage 1: Context Extraction and Sensitivity Detection}
\label{subsubsec:example2-context}

The user opens a conversation about personal stress and mental wellbeing.
The context extractor returns:

\begin{lstlisting}[style=jsonstyle, caption={Context extraction output for a
sensitive mental-health query.}, label={lst:example2-context}]
{
  "intent":          "general_question",
  "ad_receptivity":  "low",
  "purchase_signal": 0.00,
  "topics":          ["mental health", "stress", "anxiety"],
  "entities":        []
}
\end{lstlisting}

\noindent
Two signals jointly trigger suppression.  First, \texttt{purchase\_signal}
is~0.00, indicating no commercial intent whatsoever.  Second, the topics
array contains members of the hard-coded sensitivity list (\texttt{mental
health}, \texttt{anxiety}, \texttt{grief}, \texttt{financial distress},
and related terms).  The orchestrator evaluates both conditions before
dispatching a retrieval request: if either \texttt{ad\_receptivity} is
\texttt{low} \emph{and} the topic set intersects the sensitivity list, the
ad-matching stage is skipped.

\subsubsection{Stage 2: Ad Matching Skipped}
\label{subsubsec:example2-skip}

Because the sensitivity gate fires, the orchestrator forwards an empty ad
list directly to the synthesis stage without issuing any embedding or
vector-search call.  No candidates are scored, no budget gate is evaluated,
and no \texttt{<sponsored\_context>} block is assembled.  This is a stricter
outcome than the threshold-rejection path seen in
Example~\ref{subsec:example-irrelevant}: there, the pipeline searched but
found no qualifying ads; here, it does not search at all.

\subsubsection{Stage 3: Response Generation}
\label{subsubsec:example2-response}

The system prompt passed to the LLM contains only the conversational
instruction set --- no sponsored context, no tracking URL, no disclosure
obligation.  The model responds with empathetic, helpful guidance
appropriate to a mental-health topic, entirely free of product mentions.
No impression event is written and no campaign spend is recorded for this
turn.  From the user's perspective the interaction is indistinguishable from
a standard ad-free assistant; the suppression is invisible by design.
